\documentclass[../principal.tex]{subfiles}
\graphicspath{{\subfix{../imagenes/}}}

\begin{document}

\ifSubfilesClassLoaded{
    \chapter{Circuitos integrados}
}{}

\section{Chip 555}

El chip 555 es un circuito integrado que se utiliza para generar señales de temporización.

\subsection{Patitas}

El chip 555 tiene 8 patitas.

\begin{itemize}
    \item 01: GND
    \item 02: Trigger
    \item 03: Output
    \item 04: Reset
    \item 05: Control
    \item 06: Threshold
    \item 07: Discharge
    \item 08: $V_{CC}$
\end{itemize}

\subsection{Circuito astable}

La configuración astable del chip 555 nos permite oscilar entre dos estados, lo que es útil para generar señales de reloj, o de frentón para osciladores de baja frecuencia o de audio.

\subsubsection{Alimentación}

\begin{enumerate}
    \item Patita 4 conectada a $V_{CC}$
    \item Patita 8 conectada a $V_{CC}$
    \item Patita 1 conectada a GND
    \item Patita 5 conectada a GND a través de un capacitor cerámico de 100 nF
\end{enumerate}

\subsubsection{Entradas}

\begin{enumerate}
    \item Patita 2 conectada a patita 6
    \item Patita 2 conectada a tierra a través de un capacitor electrolítico de 10 \textmu F
    \item Patita 2 conectada a la patita 3 de un potenciómetro de 100 k$\Omega$
    \item Patita 7 conectada a patita 2 del mismo potenciómetro del paso anterior
    \item Patita 7 conectada a $V_{CC}$ a través de un resistor de 
\end{enumerate}


\subsubsection{Salidas}

\begin{enumerate}
    \item Patita 3 es la salida del circuito
\end{enumerate}

Recomendamos para probar la salida del circuito, conectar la patita 3 a un LED y un resistor de 220 $\Omega$ a tierra.

\subsection{Circuito monostable}



\newpage

\end{document}

